\documentclass{article}
\usepackage{graphicx} % Required for inserting images
\usepackage{amsmath}
\usepackage{physics}
\usepackage{hyperref}
\usepackage{graphicx}
\usepackage{mathrsfs}
\usepackage{siunitx}
\usepackage[section]{placeins}
\usepackage{upgreek}
\usepackage{dsfont}
\usepackage{amssymb}
\usepackage{caption}
\usepackage{subcaption}
\usepackage{multirow}
\usepackage{hhline}
\usepackage{booktabs}
\usepackage{relsize}
\usepackage{mathdots}
\usepackage{xcolor}
\usepackage{tikz}
\usepackage{pgfplots}
\usepackage{accents}
\usepackage{bbm}
\usepackage{mathtools}

\newcommand{\mathcolorbox}[2]{\mathchoice
  {\colorbox{#1}{$\displaystyle #2$}}
  {\colorbox{#1}{$\textstyle #2$}}
  {\colorbox{#1}{$\scriptstyle #2$}}
  {\colorbox{#1}{$\scriptscriptstyle #2$}}
}
\newcommand{\vc}{\boldsymbol}
\newcommand{\mt}{\boldsymbol}
\renewcommand{\d}{{\rm d}}
\renewcommand{\l}{\left}
\renewcommand{\r}{\right}
\newcommand{\ruler}{
 \begin{center}
     {\rule{8cm}{0.7pt}}
 \end{center}}
\newcommand{\trs}{^\top \!}
\newcommand{\doubleunderline}[1]{\underline{\underline{#1}}}


% Use titlesec to customize section formatting
\usepackage{titlesec}
\usepackage{titling}

% Roman numerals for sections
\renewcommand{\thesection}{\Roman{section}}

% Capital letters for subsections
\renewcommand{\thesubsection}{\Alph{subsection}}

% Center and bold section titles
\titleformat{\section}
  {\bfseries\Large}  % Format: centered, bold, large font add in \centering to make it go in the middle
  {\thesection.}{1em}{}        % Label format

\titleformat{\subsection}
  {\bfseries\normalsize}  % Format: centered, bold, normal font
  {\thesubsection.}{1em}{}          % Label format

% Optional: spacing before/after sections
\titlespacing*{\section}{0pt}{2em}{1em}
\titlespacing*{\subsection}{0pt}{1.5em}{0.5em}

\setlength{\droptitle}{-6em}
\usepackage[a4paper]{geometry}

% \newcommand{\myappendix}{
%     \newpage              % Start the appendix on a new page
%     \newpage              % Start the appendix on a new page
%     \appendix             % Switch to appendix mode
%     \renewcommand{\thesection}{A\arabic{section}} % Set appendix section format to A1, A2, etc.
%     \setcounter{section}{0} % Reset section counter to start from A1
% }
\newenvironment{myappendix}{
    \clearpage
    \appendix
    % \addcontentsline{toc}{section}{Appendix}
    \addtocontents{toc}{\protect\noindent\rule{\textwidth}{0.4pt}\par}

    \renewcommand{\thesection}{A\arabic{section}}%
    \setcounter{section}{0}%
}{
    % (optional) code to execute at the end of the appendix environment,
    % leave empty if nothing needed
}

% Adjust the \footskip to move the footer lower
\setlength{\footskip}{35pt} % Ispace between page numbers and lower text
\geometry{a4paper, top=1in, bottom=1.2in, left=1.2in, right=1.2in} % Set the margins


\usepackage[most]{tcolorbox}

% Define a custom aside box environment
\newtcolorbox{asideenv}[1]{% <-- only 1 argument, the title
  colback=gray!10,
  colframe=gray!60,
  coltitle=black,
  title={#1},
  fonttitle=\bfseries,
  boxrule=0.5pt,
  arc=3pt,
  left=6pt,
  right=6pt,
  top=6pt,
  bottom=6pt,
  enhanced,
  breakable,
}

% Define the aside environment wrapper WITHOUT numbering
\newenvironment{aside}[1]{%
  \addcontentsline{toc}{subsubsection}{#1}% Add plain title to TOC without numbers
  \begin{asideenv}{#1}%
}{%
  \end{asideenv}%
}

% for TOC putting a dot before TOC
\makeatletter
\renewcommand{\numberline}[1]{#1.\hskip0.5em}
\makeatother

\pgfplotsset{compat=1.18} % Set the compatibility level for pgfplots

\title{General Notes on FEOM}
\author{Cole Hunt}
\date{June 2025}

\begin{document}

\maketitle
This PDF contains general notes on the FEOM project, including information about the project structure, coding conventions, and other relevant details.
\section{Project Structure}
This project implements HEOM using FORTRAN, and implements basic HEOM with the standard K and L cuttoffs. The equations of motion are derived from
the Shi/Chen scaled HEOM, for which the equations of motion (without truncation) are given by
\begin{align}
    \frac{\d}{\d t}\rho_{\vc{n}}(t) = -\frac{i}{\hbar}H_s^{\cross}\rho_{\vc{n}}(t)  - \sum_k n_k\gamma_k \rho_{\vc{n}}(t) 
    + \sum_k\l( \sqrt{n_k+1}\mathcal{L}_{k-}\rho_{\vc{n}_k^+}(t) +  \sqrt{n_k}\mathcal{L}_{k+}\rho_{\vc{n}_k^-}(t)\r),
\end{align}
where $\rho_{\vc{n}}(t)$ is the full set of ADOs, $\vc{n}$ is the vector of ADO indices $\{n_k\}$, (with $\vc{n}_k^{\pm}$ denoting the same set of indices, but with $n_k\pm1$).
We use superoperator notation throughout, $H_s$ as the system Hamiltonian, the superoperators $\mathcal{L}_{k-}$ and $\mathcal{L}_{k+}$ are defined as
the superoperator 
\begin{align}
    \mathcal{L}_{k-}\cdot &= -i\sqrt{|C_k|}V^{\cross}, & \mathcal{L}_{k+}\cdot &= \frac{1}{\sqrt{|C_k|}}\l(C_k V^{L} - C_k^* V^{R}\r),
\end{align}
where $V$ is the system-bath coupling operator, $C_k$ is the coupling strength for the $k$'th bath mode, and the constants $C_k$ and $\gamma_k$ are obtained from the bath correlation function (BCF) as
\begin{align}
    C(t) = \sum_k C_k e^{-\gamma_k t}.
\end{align}

\section{Bath conventions}
In this section, we discuss the conventions used for the naming of the baths. We define the spectral density
\begin{align}
    J(\omega)&= \frac{\pi}{2}\sum_{\alpha} \frac{c_{\alpha}^2}{m_\alpha\omega_\alpha} \l[\delta(\omega - \omega_\alpha) + \delta(\omega + \omega_\alpha) \r],
\end{align}
where $c_\alpha$ is the coupling strength, $m_\alpha$ is the mass, and $\omega_\alpha$ is the frequency of the $\alpha$'th bath mode. It is assumed that $\omega_\alpha \geq 0$, such that $J(\omega)=-J(-\omega)$. 
With this definition, the equilibrium bath correlation function  (BCF), at inverse temperature $\beta$, is given as
\begin{align}
    C(t) &= \frac{\hbar}{2\pi}\int_{-\infty}^\infty d\omega J(\omega)\l(\coth\l(\frac{\beta\hbar\omega}{2}\r) \cos\omega t - i\sin\omega t\r),
\end{align}
which is done using residues.
Now to get onto specific bath distributions.

\subsection{Debye bath}
The Debye bath is defined by the spectral density
\begin{align}
    J(\omega) =\frac{\eta \gamma \omega}{\omega^2+\gamma^2},
\end{align}
where $\eta$ is the coupling strength and $\gamma$ is the cutoff frequency.
The correlation function is of exponential form, paramaterized by frequencies $\{\gamma_k\}$,
\begin{align}
    \gamma_0 &= \gamma, & \gamma_k = \omega_{n=k} = \frac{2n\pi}{\beta\hbar},
\end{align}
and prefactors $\{C_k\}$, which are given by
\begin{align}
    C_0 = \frac{\hbar\eta\gamma}{2}\l(\cot{\frac{\beta\hbar\gamma}{2}}-i\r),&
     & C_{n=k} &= -\frac{2\eta\gamma}{\beta}\frac{\omega_n}{\gamma^2-\omega_n^2}, & \forall k>0.
\end{align}

\section{Truncations}
There are a couple ways to truncate the HEOM, so it's good to write these down. 
The truncation falls into two categories, which I will call exponential term truncation and heirarchy truncation.
The exponential term truncation is the approximation of the bath correlation function (BCF) by a finite number of exponentials.
The heirarchy truncation is the removal of the dependence of the ADOs on those with a larger depth.
\subsection{Matsubara series truncation}
This trunctation is a Markovian approximation of the higher order Matsubara terms in the BCF;
\begin{align}
    e^{-\gamma_k t} &\approx \frac{1}{\gamma_k}\delta_+(t), & \forall k &> k_c,
\end{align}
where $k_c$ is the cutoff in number of exponentials in the BCF\footnote{these are ordered in increasing real part, as then the approximation is more accurate for larger $k$} and $\delta_+(t)$ is a delta function defined for $t \geq 0$.
This cutoff $k_v=M+1$ for the example of a Debye bath, wherein there is a single exponential term, and $M$ is the number of Matsubara terms retained.
This approximation is inserted into the influence functional, allowing us to do the inner integral (over the history), 
giving another term in the dynamics,
\begin{align}
    \Xi \rho_{\vc{n}}(t) &= \sum_{k>k_c} \frac{1}{\gamma_k} \mathcal{L}_{k-}\mathcal{L}_{k+} \rho_{\vc{n}}(t).
\end{align}
For the example of a Debye bath, we can do the sum formally. Due to all of the $C_k$ for $k>0$ being real, we can write the superoperator as
\begin{align}
     \Xi_{\rm debye}  &= -i\sum_{k>k_c} \frac{1}{\gamma_k}C_k \l(V^{\cross}\r)^2= \frac{2i\eta\gamma}{\beta}\sum_{n>k_c} \frac{1}{\gamma^2-\omega_n^2}\l(V^{\cross}\r)^2
\end{align}
Then we use the fact that
\begin{align}
    \cot z &= \frac{1}{z} + \sum_{n=1}^\infty \frac{2z}{z^2-n^2\pi^2} \nonumber \\[6pt]
    \therefore \sum_{n=1}^\infty \frac{1}{\gamma^2-\omega_n^2} &= \frac{\beta\hbar}{4\gamma}\l(\cot{\frac{\beta\hbar\gamma}{2}}-\frac{2}{\beta\hbar\gamma}\r),
\end{align}
to calculate the sum, giving us
\begin{align}
    \Xi_{\rm debye} &= i\eta\l[-\frac{\hbar}{2}\l(\cot{\frac{\beta\hbar\gamma}{2}}-\frac{2}{\beta\hbar\gamma}\r)- \frac{2\gamma}{\beta}\sum_{n=1}^{k_c} \frac{1}{\gamma^2-\omega_n^2}\r]\l(V^{\cross}\r)^2.
\end{align}
\end{document}



